% Introduction condensée (FR). Traduction de parts/01_introduction.tex, style
% CamemBERT-bio (TALN 2023). Formulations reprises du résumé officiel de la thèse.

\partepigraph{``The biggest lesson that can be read from 70 years of AI research is
that general methods that leverage computation are ultimately the most effective, and
by a large margin.''}{Richard S. Sutton, \textit{The Bitter Lesson}}

\noindent Comment apprendre à une intelligence artificielle à lire des comptes-rendus
hospitaliers quand ces comptes-rendus ne peuvent pas servir à l'entraîner ? Ils
seraient pourtant les meilleurs exemples, car ils contiennent les mots, les
raccourcis et les habitudes cliniques que le système doit apprendre. Mais ils
décrivent des patients réels : ils restent donc à l'intérieur de l'hôpital, et un
modèle entraîné sur ces données pourrait devoir y rester aussi. Cette thèse étudie
si du texte biomédical public peut être utilisé à la place.

La question est ancienne. En 1959, Robert Ledley et Lee Lusted donnent au nouveau
domaine de l'intelligence artificielle un objectif médical, la création d'un ordinateur capable
de pondérer les signes d'un patient vers une maladie probable~\citep{ledley_lusted_1959}.
Leur question est aussi celle de cette thèse : d'où vient la connaissance qui fait
fonctionner un tel système, et peut-on en fournir assez à une machine pour qu'elle
soit utile aux soins ? Les premiers systèmes médicaux répondaient à partir de leurs
propres dossiers. Le système de \citet{dedombal_1972} pour l'abdomen aigu atteint
91,8\% d'exactitude diagnostique sur 304 patients à Leeds en 1972. Repris dans un
essai multicentrique sur huit centres et 16\,737 patients, il n'y est plus évalué
seul, mais comme une aide : l'exactitude initiale des médecins passe de 45,6\% avant
son installation à 65,3\% avec lui~\citep{adams_1986}. L'essai est un succès, et ce
qui n'a pas voyagé, c'est l'ambition de 1972 : ailleurs, le programme a été placé à
côté du clinicien plutôt que devant lui, car sa connaissance était une table de
fréquences venue d'un seul hôpital. Un système qui apprend dépend des données à partir desquelles
sa connaissance a été construite ; un résultat mesuré dans un cadre peut donc ne pas
se vérifier dans un autre.

Dans les années 1970, l'IA médicale se tourne vers les \textit{systèmes experts},
qui représentent les connaissances des spécialistes sous forme de règles explicites.
À Stanford, MYCIN recommandait des antibiotiques à l'aide de règles
SI--ALORS~\citep{mycin_1975}. Lors d'une évaluation en aveugle, huit spécialistes des
maladies infectieuses jugèrent ses recommandations acceptables dans 65\% des cas,
contre 42,5\% à 62,5\% pour les cinq médecins enseignants~\citep{mycin_eval_1979}. Malgré ces
résultats, MYCIN ne fut jamais utilisé en pratique clinique. Ses règles devaient être
rédigées et vérifiées une à une par des spécialistes, puis mises à jour à mesure que
les connaissances médicales évoluaient. Maintenir une telle base était coûteux, et son
exhaustivité ne pouvait jamais être garantie. Les systèmes experts fonctionnaient donc
sur des tâches étroites, mais leur connaissance construite à la main passait
difficilement à l'échelle.

L'apprentissage automatique a changé cette manière d'acquérir la connaissance.
Plutôt que de recevoir les règles, les systèmes apprennent quoi faire à
partir de grands ensembles d'exemples, et les gains ont été importants : dès 2016, un
réseau entraîné sur 128\,175 photographies de rétine les lisait pour la rétinopathie
diabétique avec une exactitude proche de celle des ophtalmologues~\citep{gulshan_2016}.
Dans la littérature biomédicale, le nombre d'articles sur les systèmes experts
culmine vers 1991, tandis que les publications consacrées à l'apprentissage profond,
encore rares avant 2015, dépassent ensuite les vingt mille par an. Cette évolution
pose le problème au cœur de cette thèse : l'apprentissage automatique exige beaucoup
de données adaptées à la tâche, alors que les textes cliniques sont rarement
accessibles.

La plupart des traitements automatiques des langues modernes commencent par
pré-entraîner un modèle de langue sur de grandes collections de texte brut, avant de
l'adapter à une tâche. Dans le cadre clinique, cette voie est difficile à suivre. Les
comptes-rendus cliniques sont des données personnelles et ne peuvent être partagés :
la France n'a donc pas de corpus public sur lequel les pré-entraîner. Un modèle
entraîné à l'intérieur d'un hôpital sur ses propres dossiers hérite de la même
contrainte et ne peut être diffusé aux autres. Un modèle clinique partageable doit
être construit à partir de texte public. Ce manque de données limite aussi
l'évaluation. Sans dossiers cliniques accessibles, il est aussi difficile de mesurer
les capacités cliniques des modèles que de les entraîner. Le matériel utilisé pour
les évaluer est donc reconstruit à partir de sources publiques, puis examiné par des
cliniciens de CHU français.

Cette thèse pose une seule question : des données publiques peuvent-elles soutenir
des modèles de langue pour l'extraction d'information clinique en français, sans
entraînement sur des dossiers cliniques sensibles ? Le manuscrit y répond à travers
trois questions.

\begin{itemize}
  \item Pouvons-nous constituer un corpus biomédical français public utile, et son contenu et sa qualité changent-ils ce que le modèle apprend ?
  \item Étant donné un tel corpus, comment pré-entraîner un modèle de langue en combinant l'adaptation d'un modèle existant, le choix de l'objectif d'entraînement et l'ajout de connaissances extérieures au texte ?
  \item Comment adapter ces modèles à des tâches d'extraction clinique quand les données annotées sont rares, par l'architecture du modèle et par des données synthétiques ?
\end{itemize}

\noindent Cette thèse apporte sept contributions.

\begin{itemize}
  \item \emph{biomed-fr}, un corpus biomédical français public de 413 millions de mots, utilisé pour entraîner CamemBERT-bio. Le modèle adapté repère les termes biomédicaux d'un texte français, comme les médicaments et les maladies, plus précisément que le modèle généraliste dont il part, avec un gain moyen de 2,54 points et pour une fraction du coût d'un entraînement à partir de zéro~\citep{touchent2023camembertbio}.
  \item \emph{Biomed-Enriched}, un étiquetage de la littérature biomédicale ouverte (PubMed Central) par type de contenu et par qualité. Il identifie environ deux millions de paragraphes de cas cliniques et améliore un modèle dont on poursuit l'entraînement sur ces paragraphes~\citep{touchent-etal-2026-biomed}.
  \item \emph{biomed-fr-v2}, une amélioration du corpus biomédical français avec dix milliards de tokens, et une étude des signaux de qualité qui aident le plus. Certains filtres courants se révèlent dégrader les performances.
  \item \emph{ModernCamemBERT-bio}, le premier modèle français biomédical et clinique à contexte long, construit avec une méthode de pré-entraînement qui améliore l'adaptation au domaine.
  \item \emph{OntoBook}, une approche et un jeu de données qui transforment une base de connaissances médicale en texte d'entraînement et améliorent le codage médical français.
  \item \emph{MC-bio-gliner} et \emph{MC-bio-embed}, deux variantes de ModernCamemBERT-bio, l'une qui extrait de l'information des comptes-rendus cliniques, l'autre qui apparie des termes à des codes médicaux, toutes deux à partir de peu de données annotées.
  \item Une méthode pour générer des comptes-rendus cliniques synthétiques en réécrivant du texte biomédical public dans le style hospitalier, afin d'entraîner des extracteurs sans accès à de vraies données hospitalières.
\end{itemize}

\noindent Les parties suivent ces questions. La première demande ce qu'est un texte biomédical
et à quel point ses formes publiques sont éloignées des comptes-rendus cliniques qui
nous intéressent. La deuxième construit et étudie un corpus biomédical français
public. La troisième pré-entraîne des modèles de langue dessus, par adaptation de
modèle, choix de l'objectif et enrichissement par la connaissance. La dernière partie
adapte ces modèles à des tâches d'extraction clinique sous faible quantité de données
annotées. La conclusion ouvre alors sur une question que ce travail laisse derrière
lui : une fois ces modèles construits, la manière dont ils se comportent lorsqu'on les
utilise réellement devient un problème en soi.
